% De Lege Agraria I (de-lege-agraria-1)
% Source: https://raw.githubusercontent.com/PerseusDL/canonical-latinLit/master/data/phi0474/phi011/phi0474.phi011.perseus-lat2.xml
% Fetched by scripts/fetch_latin.py

\ciceroSection{1f} imberba iuventute

\ciceroSection{2f} Capuam colonis deductis occupabunt, Atellam praesidio communient, Nuceriam, Cumas multitudine suorum obtinebunt, cetera oppida praesidiis devincient.

\ciceroSection{3f} Venibit igitur sub praecone tota Propontis atque Hellespontus, addicetur omnis ora Lyciorum atque Cilicum, Mysia et Phrygia eidem condicioni legique parebunt.

\ciceroSection{4f} praedam , manubias, sectionem, castra denique Cn. Pompei sedente imperatore x viri vendent.

\ciceroSection{1} * * * quae res aperte petebatur, ea nunc occulte cuniculis oppugnatur. dicent enim x viri, id quod et dicitur a multis et saepe dictum est, post eosdem consules regis Alexandri testamento regnum illud populi Romani esse factum. dabitis igitur Alexandream clam petentibus eis quibus apertissime pugnantibus restitistis? haec , per deos immortalis! utrum esse vobis consilia siccorum an vinolentorum somnia, et utrum cogitata sapientium an optata furiosorum videntur?

\ciceroSection{2} videte nunc proximo capite ut impurus helluo turbet rem publicam, ut a maioribus nostris possessiones relictas disperdat ac dissipet, ut sit non minus in populi Romani patrimonio nepos quam in suo. perscribit in sua lege vectigalia quae x viri vendant, hoc est, proscribit auctionem publicorum bonorum. agros emi volt qui dividantur; quaerit pecuniam. videlicet excogitabit aliquid atque adferet. nam superioribus capitibus dignitas populi Romani violabatur, nomen imperi in commune odium orbis terrae vocabatur, urbes pacatae, agri sociorum, regum status x viris donabantur; nunc praesens pecunia, certa, numerata quaeritur.

\ciceroSection{3} exspecto quid tribunus plebis vigilans et acutus excogitet. 'Veneat,' inquit, 'silva Scantia.' Vtrum tandem hanc silvam in relictis possessionibus, an in censorum pascuis invenisti? si quid est quod indagaris, inveneris, ex tenebris erueris, quamquam iniquum est, tamen consume sane, quod commodum est, quoniam quidem tu attulisti; silvam vero tu Scantiam vendas nobis consulibus atque hoc senatu? tu ullum vectigal attingas, tu populo Romano subsidia belli, tu ornamenta pacis eripias? tum vero hoc me inertiorem consulem iudicabo quam illos fortissimos viros qui apud maiores nostros fuerunt, quod, quae vectigalia illis consulibus populo Romano parta sunt, ea me consule ne retineri quidem potuisse iudicabuntur.

\ciceroSection{4} vendit Italiae possessiones ex ordine omnis. sane est in eo diligens; nullam enim praetermittit. persequitur in tabulis censoriis totam Siciliam; nullum aedificium, nullos agros relinquit. Audistis auctionem populi Romani proscriptam a tribuno plebis, constitutam in mensem Ianuarium, et, credo, non dubitatis quin idcirco haec aerari causa non vendiderint ei qui armis et virtute pepererunt, ut esset quod nos largitionis causa venderemus.

\ciceroSection{5} videte nunc quo adfectent iter apertius quam antea. nam superiore parte legis quem ad modum Pompeium oppugnarent, a me indicati sunt; nunc iam se ipsi indicabunt. iubent venire agros Attalensium atque Olympenorum quos populo Romano P. Servili, fortissimi viri, victoria adiunxit, deinde agros in Macedonia regios qui partim T. Flaminini, partim L. Pauli qui Persen vicit virtute parti sunt, deinde agrum optimum et fructuosissimum Corinthium qui L. Mummi imperio ac felicitate ad vectigalia populi Romani adiunctus est, post autem agros in Hispania apud Carthaginem novam duorum Scipionum eximia virtute possessos; tum vero ipsam veterem Carthaginem vendunt quam P. Africanus nudatam tectis ac moenibus sive ad notandam Carthaginiensium calamitatem, sive ad testificandam nostram victoriam, sive oblata aliqua religione ad aeternam hominum memoriam consecravit.

\ciceroSection{6} his insignibus atque infulis imperi venditis quibus ornatam nobis maiores nostri rem publicam tradiderunt, iubent eos agros venire quos rex Mithridates in Paphlagonia, Ponto Cappadociaque possederit. num obscure videntur prope hasta praeconis insectari Cn. Pompei exercitum qui venire iubeant eos ipsos agros in quibus ille etiam nunc bellum gerat atque versetur?

\ciceroSection{7} hoc vero cuius modi est, quod eius auctionis quam constituunt locum sibi nullum definiunt? nam x viris quibus in locis ipsis videatur vendendi potestas lege permittitur. censoribus vectigalia locare nisi in conspectu populi Romani non licet; his vendere vel in ultimis terris licebit? at hoc etiam nequissimi homines consumptis patrimoniis faciunt ut in atriis auctionariis potius quam in triviis aut in compitis auctionentur; hic permittit sua lege x viris ut in quibus commodum sit tenebris, ut in qua velint solitudine, bona populi Romani possint divendere.

\ciceroSection{8} iam illa omnibus in provinciis, regnis, liberis populis quam acerba, quam formidolosa, quam quaestuosa concursatio x viralis futura sit, non videtis? hereditatum obeundarum causa quibus vos legationes dedistis, qui et privati et privatum ad negotium exierunt non maximis opibus neque summa auctoritate praediti, tamen auditis profecto quam graves eorum adventus sociis nostris esse soleant.

\ciceroSection{9} quam ob rem quid putatis impendere hac lege omnibus gentibus terroris et mali, cum immittantur in orbem terrarum x viri summo cum imperio, summa cum avaritia infinitaque omnium rerum cupiditate? quorum cum adventus graves, cum fasces formidolosi, tum vero iudicium ac potestas erit non ferenda; licebit enim quod videbitur publicum iudicare, quod iudicarint vendere. etiam illud quod homines sancti non facient, ut pecuniam accipiant ne vendant, tamen id eis ipsum per legem licebit. hinc vos quas spoliationes, quas pactiones, quam denique in omnibus locis nundinationem iuris ac fortunarum fore putatis?

\ciceroSection{10} etenim , quod superiore parte legis praefinitum fuit, ' Svlla et Pompeio consvlibvs ,' id rursus liberum infinitumque fecerunt. iubet enim eosdem x viros omnibus agris publicis pergrande vectigal imponere, ut idem possint et liberare agros quos commodum sit et quos ipsis libeat publicare. quo in iudicio perspici non potest utrum severitas acerbior an benignitas quaestuosior sit futura. sunt tamen in tota lege exceptiones duae non tam iniquae quam suspiciosae. excipit enim in vectigali imponendo agrum Recentoricum Siciliensem, in vendendis agris eos agros de quibus cautum sit foedere. hi sunt in Africa, qui ab Hiempsale possidentur.

\ciceroSection{11} hic quaero, si Hiempsali satis est cautum foedere et Recentoricus ager privatus est, quid attinuerit excipi; sin et foedus illud habet aliquam dubitationem et ager Recentoricus dicitur non numquam esse publicus, quem putet existimaturum duas causas in orbe terrarum repertas quibus gratis parceret. num quisnam tam abstrusus usquam nummus videtur quem non architecti huiusce legis olfecerint? provincias , civitates liberas, socios, amicos, reges denique exhauriunt, admovent manus vectigalibus populi Romani.

\ciceroSection{12} non est satis. audite , audite vos qui amplissimo populi senatusque iudicio exercitus habuistis et bella gessistis: quod ad quemque pervenerit ex praeda, ex manubiis, ex auro coronario, quod neque consumptum in monumento neque in aerarium relatum sit, id ad x viros referri iubet! hoc capite multa sperant; in omnis imperatores heredesque eorum quaestionem suo iudicio comparant, sed maximam pecuniam se a Fausto ablaturos arbitrantur. quam causam suscipere iurati iudices noluerunt, hanc isti x viri susceperunt: idcirco a iudicibus fortasse praetermissam esse arbitrantur quod sit ipsis reservata.

\ciceroSection{13} deinde etiam in reliquum tempus diligentissime sancit ut, quod quisque imperator habeat pecuniae, protinus ad x viros deferat. hic tamen excipit Pompeium simillime, ut mihi videtur, atque ut illa lege qua peregrini Roma eiciuntur Glaucippus excipitur. non enim hac exceptione unus adficitur beneficio, sed unus privatur iniuria. sed cui manubias remittit, in huius vectigalia invadit. iubet enim pecunia, si qua post nos consules ex novis vectigalibus recipiatur, hac uti x viros. quasi vero non intellegamus haec eos vectigalia quae Cn. Pompeius adiunxerit vendere cogitare.

\ciceroSection{14} videtis iam, patres conscripti, omnibus rebus et modis constructam et coacervatam pecuniam x viralem. minuetur huius pecuniae invidia; consumetur enim in agrorum emptionibus. optime . quis ergo emet agros istos? idem x viri; tu, Rulle,—missos enim facio ceteros—emes quod voles, vendes quod voles; utrumque horum facies quanti voles. cavet enim vir optimus ne emat ab invito. quasi vero non intellegamus ab invito emere iniuriosum esse, ab non invito quaestuosum. quantum tibi agri vendet, ut alios omittam, socer tuus, et, si ego eius aequitatem animi probe novi, vendet non invitus? facient idem ceteri libenter, ut possessionis invidiam pecunia commutent, accipiant quod cupiunt, dent quod retinere vix possunt.

\ciceroSection{15} nunc perspicite omnium rerum infinitam atque intolerandam licentiam. pecunia coacta est ad agros emendos; ei porro ab invitis non ementur. si consenserint possessores non vendere, quid futurum est? referetur pecunia? non licet. exigetur ? vetat . verum esto; nihil est quod non emi possit, si tantum des quantum velit venditor. spoliemus orbem terrarum, vendamus vectigalia, effundamus aerarium, ut locupletatis aut invidiae aut pestilentiae possessoribus agri tamen emantur.

\ciceroSection{16} quid tum? quae erit in istos agros deductio, quae totius rei ratio atque descriptio? ' deducentur ,' inquit, 'coloniae.' quot ? quorum hominum? in quae loca? quis enim non videt in coloniis esse haec omnia consideranda? tibi nos, Rulle, et istis tuis harum omnium rerum machinatoribus totam Italiam inermem tradituros existimasti, quam praesidiis confirmaretis, coloniis occuparetis, omnibus vinclis devinctam et constrictam teneretis? Vbi enim cavetur ne in Ianiculo coloniam constituatis, ne urbem hanc urbe alia premere atque urgere possitis? ' non faciemus,' inquit. primum nescio, deinde timeo, postremo non committam ut vestro beneficio potius quam nostro consilio salvi esse possimus.

\ciceroSection{17} quod vero totam Italiam vestris coloniis complere voluistis, id cuius modi esset neminemne nostrum intellecturum existimavistis? scriptum est enim: ' Qvae in mvnicipia qvasqve in colonias xviri velint, dedvcant colonos qvos velint et eis agros adsignent qvibvs in locis velint ,' ut, cum totam Italiam militibus suis occuparint, nobis non modo dignitatis retinendae, verum ne libertatis quidem recuperandae spes relinquatur. atque haec a me suspicionibus et coniectura coarguuntur.

\ciceroSection{18} iam omnis omnium tolletur error, iam aperte ostendent sibi nomen huius rei publicae, sedem urbis atque imperi, denique hoc templum Iovis optimi maximi atque hanc arcem omnium gentium displicere. Capuam deduci colonos volunt, illam urbem huic urbi rursus opponere, illuc opes suas deferre et imperi nomen transferre cogitant. qui locus propter ubertatem agrorum abundantiamque rerum omnium superbiam et crudelitatem genuisse dicitur, ibi nostri coloni delecti ad omne facinus a x viris conlocabuntur, et, credo, qua in urbe homines in vetere dignitate fortunaque nati copiam rerum moderate ferre non potuerunt, in ea isti vestri satellites modeste insolentiam suam continebunt.

\ciceroSection{19} maiores nostri Capua magistratus, senatum, consilium commune, omnia denique insignia rei publicae sustulerunt, neque aliud quicquam in urbe nisi inane nomen Capuae reliquerunt, non crudelitate—quid enim illis fuit clementius qui etiam externis hostibus victis sua saepissime reddiderunt?—sed consilio, quod videbant, si quod rei publicae vestigium illis moenibus contineretur, urbem ipsam imperio domicilium praebere posse; vos haec, nisi evertere rem publicam cuperetis ac vobis novam dominationem comparare, credo, quam perniciosa essent non videretis.

\ciceroSection{20} quid enim cavendum est in coloniis deducendis? si luxuries, Hannibalem ipsum Capua corrupit, si superbia, nata inibi esse haec ex Campanorum fastidio videtur, si praesidium, non praeponitur huic urbi ista colonia, sed opponitur. at quem ad modum armatur, di immortales! nam bello Punico quicquid potuit Capua, potuit ipsa per sese; nunc omnes urbes quae circum Capuam sunt a colonis per eosdem x viros occupabuntur; hanc enim ob causam permittit ipsa lex, in omnia quae velint oppida colonos ut x viri deducant quos velint. atque his colonis agrum Campanum et Stellatem campum dividi iubet.

\ciceroSection{21} non queror deminutionem vectigalium, non flagitium huius iacturae atque damni, praetermitto illa quae nemo est quin gravissime et verissime conqueri possit, nos caput patrimoni publici, pulcherrimam populi Romani possessionem, subsidium annonae, horreum belli, sub signo claustrisque rei publicae positum vectigal servare non potuisse, eum denique nos agrum P. Rullo concessisse, qui ager ipse per sese et Sullanae dominationi et Gracchorum largitioni restitisset; non dico solum hoc in re publica vectigal esse quod amissis aliis remaneat, intermissis non conquiescat, in pace niteat, in bello non obsolescat, militem sustentet, hostem non pertimescat; praetermitto omnem hanc orationem et contioni reservo; de periculo salutis ac libertatis loquor.

\ciceroSection{22} quid enim existimatis integrum vobis in re publica fore aut in vestra libertate ac dignitate retinenda, cum Rullus atque ei quos multo magis quam Rullum timetis cum omni egentium atque improborum manu, cum omnibus copiis, cum omni argento et auro Capuam et urbis circa Capuam occuparint? his ego rebus, patres conscripti, resistam vehementer atque acriter neque patiar homines ea me consule expromere quae contra rem publicam iam diu cogitarunt.

\ciceroSection{23} Errastis, Rulle, vehementer et tu et non nulli conlegae tui qui sperastis vos contra consulem veritate, non ostentatione popularem posse in evertenda re publica populares existimari. lacesso vos, in contionem voco, populo Romano disceptatore uti volo. etenim , ut circumspiciamus omnia quae populo grata atque iucunda sunt, nihil tam populare quam pacem, quam concordiam, quam otium reperiemus. sollicitam mihi civitatem suspicione, suspensam metu, perturbatam vestris legibus et contionibus et deductionibus tradidistis; spem improbis ostendistis, timorem bonis iniecistis, fidem de foro, dignitatem de re publica sustulistis.

\ciceroSection{24} hoc motu atque hac perturbatione animorum atque rerum cum populo Romano vox et auctoritas consulis repente in tantis tenebris inluxerit, cum ostenderit nihil esse metuendum, nullum exercitum, nullam manum, nullas colonias, nullam venditionem vectigalium, nullum imperium novum, nullum regnum x virale, nullam alteram Romam neque aliam sedem imperi nobis consulibus futuram summamque tranquillitatem pacis atque oti, verendum, credo, nobis erit ne vestra ista praeclara lex agraria magis popularis esse videatur.

\ciceroSection{25} Cum vero scelera consiliorum vestrorum fraudemque legis et insidias quae ipsi populo Romano a popularibus tribunis plebis fiant ostendero, pertimescam, credo, ne mihi non liceat contra vos in contione consistere, praesertim cum mihi deliberatum et constitutum sit ita gerere consulatum quo uno modo geri graviter et libere potest, ut neque provinciam neque honorem neque ornamentum aliquod aut commodum neque rem ullam quae a tribuno plebis impediri possit appetiturus sim.

\ciceroSection{26} dicit frequentissimo senatu consul Kalendis Ianuariis sese, si status hic rei publicae maneat neque aliquod negotium exstiterit quod honeste subterfugere non possit, in provinciam non iturum. sic me in hoc magistratu geram, patres conscripti, ut possim tribunum plebis rei publicae iratum coercere, mihi iratum contemnere. quam ob rem, per deos immortalis! conligite vos, tribuni plebis, deserite eos a quibus, nisi prospicitis, brevi tempore deseremini, conspirate nobiscum, consentite cum bonis, communem rem publicam communi studio atque amore defendite. multa sunt occulta rei publicae volnera, multa nefariorum civium perniciosa consilia; nullum externum periculum est, non rex, non gens ulla, non natio pertimescenda est; inclusum malum, intestinum ac domesticum est. huic pro se quisque nostrum mederi atque hoc omnes sanare velle debemus.

\ciceroSection{27} erratis , si senatum probare ea quae dicuntur a me putatis, populum autem esse in alia voluntate. omnes qui se incolumis volent sequentur auctoritatem consulis soluti a cupiditatibus, liberi a delictis, cauti in periculis, non timidi in contentionibus. quod si qui vestrum spe ducitur se posse turbulenta ratione honori velificari suo, primum me consule id sperare desistat, deinde habeat me ipsum sibi documento, quem equestri ortum loco consulem videt, quae vitae via facillime viros bonos ad honorem dignitatemque perducat. quod si vos vestrum mihi studium, patres conscripti, ad communem dignitatem defendendam profitemini, perficiam profecto, id quod maxime res publica desiderat, ut huius ordinis auctoritas, quae apud maiores nostros fuit, eadem nunc longo intervallo rei publicae restituta esse videatur.
